\section{Glossary: Reader-Facing Names and Repository Labels}
\label{app:glossary}

The paper uses reader-facing descriptive names in body text and introduces the
matching repository label once, in parentheses, on first use. This appendix
collects the full mapping for readers who encounter a label in the codebase or
in committed artifact paths.

{\small
\setlength{\LTpre}{6pt}
\setlength{\LTpost}{6pt}
\setlength{\tabcolsep}{4pt}
\begin{longtable}{@{}>{\raggedright\arraybackslash}p{0.30\linewidth} >{\raggedright\arraybackslash}p{0.26\linewidth} >{\raggedright\arraybackslash}p{0.40\linewidth}@{}}
\caption[Glossary: descriptive names and repository labels.]{Glossary:
descriptive names and repository labels.}\label{tab:glossary}\\
\toprule
Reader-facing name & Repo label & Meaning \\
\midrule
\endfirsthead
\toprule
Reader-facing name & Repo label & Meaning \\
\midrule
\endhead
prompt-regression guard & Phase~A & Pre-scoring check that the extractor's
  prompt has not regressed against locked reference outputs. \\
\addlinespace
noisy extractor-quality gate & Phase~3 & Component-quality gate that scores
  local extractor outputs and licenses downstream policy comparison. \\
\addlinespace
noisy same-stream policy comparison & Phase~4 & The main noisy-pipeline
  policy comparison that runs only after the extractor-quality gate passes. \\
\addlinespace
cardinality-repair follow-up & Phase~Y & The LongMemEval follow-up that
  isolates answer-time readout cardinality as the mechanism behind the
  transfer-probe loss. \\
\addlinespace
LongMemEval transfer probe & X.4 & The controlled-pilot external transfer of
  the same-stream protocol to LongMemEval v1. \\
\addlinespace
preregistered-gate condition fired & Bucket~A & A planned gate or thesis
  condition fired during a preregistered comparison. \\
\addlinespace
mixed/partial support & Bucket~B & A completed result that is mixed or
  partial; rows split between mechanism-survival, partial survival, and
  bounded nulls. \\
\addlinespace
abort / unsupported claim & Bucket~D & A preregistered abort or surprise-lane
  boundary that prevents the stronger claim. \\
\addlinespace
lookup diagnostic & QR-canon & The canonical-id PFLC instance: a small
  committed diagnostic table over per-question canonical ids. \\
\addlinespace
resolution audit & CQR & The earlier canonical-id/query-resolution replay
  audit whose set-membership metric and alias function the lookup diagnostic
  reuses. \\
\addlinespace
Mem0-style write-only baseline & \splitcode{Mem0\allowbreak Lite} (repo);
  \splitcode{Mem0\allowbreak Write\allowbreak Policy\allowbreak Lite} (paper)
  & Narrow ADD/UPDATE/NOOP write-time baseline; \emph{not} an end-to-end
  reproduction of the production Mem0 system. Repo identifier kept stable
  for artifact compatibility; paper-facing label changed to remove brand
  collision. \\
\addlinespace
oracle-mode dated-evidence repair & \splitcode{CQDated\allowbreak Contestation}
  (macro \code{\textbackslash cqdated}) & Oracle-mode repair variant that
  prefers dated evidence on temporal-skew rows; used in the adversarial
  upstream-noise follow-up. \\
\addlinespace
CQ plural-pending readout variant & \splitcode{CQPending\allowbreak
  MultiEvidence} (macro \code{\textbackslash cqmulti}) & CQ-side follow-up
  that exposes all eligible same-slot pending candidates at answer time;
  the repair half of the cardinality-repair follow-up. \\
\addlinespace
symmetric capped-Reflection control & \splitcode{Reflection\allowbreak
  EagerWrite\allowbreak Cardinality\allowbreak Capped} (macro
  \code{\textbackslash refcapped}) & Reflection-side control with answer
  readout capped to one durable candidate id; the symmetric mechanism
  falsifier in the cardinality-repair follow-up. \\
\bottomrule
\end{longtable}
}

The 7B and 32B labels in body text refer to local extractor model cells used
by the noisy pipeline (Qwen2.5-7B-Instruct and Qwen2.5-32B-Instruct, both
quantized), not memory-policy sizes.
